% =============================================================================
% CAPÍTULO 3: Marco Teórico / Estado del Arte
% =============================================================================
\chapterUnnumbered{Capítulo 3: Marco Teórico / Estado del Arte}

El desarrollo de soluciones basadas en Inteligencia Artificial (IA) para la síntesis de datos clínicos no puede entenderse como un esfuerzo puramente algorítmico; se trata de una convergencia multidisciplinar que entrelaza la criptografía, la estadística bayesiana, el aprendizaje profundo y el derecho tecnológico. En este capítulo se establece el marco teórico y se revisa el estado del arte de las tecnologías e investigaciones que fundamentan el presente Trabajo de Fin de Máster.

El capítulo se ha estructurado en cuatro grandes bloques lógicos. En la primera sección se analiza el paradigma de la privacidad de datos, disecando el marco normativo europeo y los criterios legales que definen la anonimización. La segunda sección traza un recorrido evolutivo por las arquitecturas de Deep Learning generativo clásico, analizando rigurosamente las Redes Generativas Antagónicas (GANs) y los Autoencoders Variacionales (VAEs). La tercera sección aborda la vanguardia tecnológica: los Modelos Probabilísticos de Difusión (Diffusion Models) y su reciente adaptación al ecosistema de los datos tabulares. Finalmente, la cuarta sección consolida los protocolos formales de evaluación matemática requeridos para cuantificar la fidelidad, la utilidad y la seguridad de los gemelos digitales generados.

\setcounter{section}{0}
\section{El Paradigma de la Privacidad de Datos en el Sector Clínico}

La proliferación de los sistemas de información hospitalarios ha generado vastos repositorios de datos masivos. Sin embargo, el historial clínico de un paciente no es meramente un conjunto de variables estadísticas; constituye una proyección digital íntima de su salud, genética y vulnerabilidades. La fuga, explotación o reidentificación indebida de estos datos puede originar consecuencias devastadoras, desde la discriminación en la concesión de seguros de salud y oportunidades laborales, hasta el fraude de identidad y la extorsión.

\subsection{El Reglamento General de Protección de Datos (GDPR)}

A nivel europeo, la protección de la información ciudadana está regida por el Reglamento (UE) 2016/679, universalmente conocido como el Reglamento General de Protección de Datos \parencite{gdpr2016}. El GDPR establece un marco punitivo y regulatorio sumamente estricto para el manejo de Información de Identificación Personal (PII, por sus siglas en inglés). En el contexto médico, los historiales clínicos son catalogados como ``categorías especiales de datos personales'' (Artículo 9 del GDPR), lo cual prohíbe explícitamente su tratamiento y distribución generalizada salvo en circunstancias de consentimiento explícito o interés público ineludible.

Esta severidad regulatoria ha paralizado intrínsecamente la colaboración en el ecosistema científico. Las instituciones de investigación se topan con el denominado \textit{Time-to-Data}, un periodo burocrático de meses requerido para tramitar Acuerdos de Procesamiento de Datos, superar Comités Éticos y evaluar el Impacto en la Privacidad (DPIA). Para mitigar este cuello de botella, el propio GDPR proporciona una vía de escape: si los datos son rigurosamente ``anónimos'', dejan de estar bajo la jurisdicción del reglamento, permitiendo su libre procesamiento.

\subsection{Antecedentes Históricos: Del k-anonimato a la Síntesis Generativa}

Antes de que la Inteligencia Artificial Generativa se afianzara como la solución técnica preferente, la comunidad científica desplegó una amplia batería de técnicas estadísticas clásicas para intentar cumplir con los requisitos de anonimización. Un análisis crítico de sus limitaciones es indispensable para contextualizar la necesidad del enfoque adoptado en el presente trabajo.

\textbf{El Modelo de k-anonimato.} La primera propuesta sistemática y formal fue el \textit{k-anonimato}, introducido por \textcite{sweeney2002kanon}. Bajo este paradigma, un conjunto de datos satisface el k-anonimato si cada registro es indistinguible de al menos $k-1$ registros adicionales con respecto a los denominados \textit{quasi-identificadores} (atributos como la edad, el código postal o el género que, individualmente, son inofensivos pero que combinados pueden singularizar a un individuo). Formalmente, para un conjunto de datos $D$ y un subconjunto de atributos $QI$, la condición de k-anonimato exige que:
$$ \forall r \in D, \; |\{ s \in D \mid s[QI] = r[QI] \}| \geq k $$
Para alcanzar esta propiedad, se aplican técnicas de \textit{supresión} (eliminación de registros que no pueden generalizarse) o \textit{generalización} (sustitución de valores precisos por rangos, como reemplazar la edad ``34'' por el intervalo ``[30-40]''). A pesar de su elegancia conceptual, el k-anonimato adolece de dos vulnerabilidades críticas: los \textbf{ataques de homogeneidad}, donde todos los registros del grupo $k$ comparten el mismo valor sensible (haciendo trivial la inferencia), y los \textbf{ataques de conocimiento previo}, donde el adversario conoce datos contextuales sobre el individuo objetivo que le permiten reidentificarlo incluso dentro de un grupo de $k$ registros aparentemente equivalentes.

\textbf{La l-diversidad como extensión.} Para subsanar la vulnerabilidad ante los ataques de homogeneidad, \textcite{machanavajjhala2007ldiv} propusieron el modelo de \textit{l-diversidad}. Este criterio extiende el k-anonimato exigiendo que cada grupo de $k$ registros equivalentes contenga al menos $l$ valores ``bien representados'' del atributo sensible. A pesar de que la l-diversidad mitiga los ataques de homogeneidad, introducen sus propias debilidades: el \textbf{ataque de similitud}, operativo cuando los $l$ valores distintos del atributo sensible son semánticamente próximos (por ejemplo, tres tipos diferentes de enfermedad renal seguirían revelando información médica crítica); y la dificultad computacional de alcanzar la l-diversidad en conjuntos de datos de alta dimensionalidad sin destruir su integridad estadística.

\textbf{El colapso empírico de la anonimización estadística.} La limitación más devastadora de todos los enfoques clásicos no fue teórica, sino empírica. En un experimento seminal, \textcite{narayanan2008netflix} demostraron que el conjunto de datos de valoraciones de películas de Netflix, publicado como supuestamente anónimo, podía ser reidentificado de forma automatizada cruzándolo con valoraciones públicas disponibles en la plataforma IMDb. Aprovechando únicamente un subconjunto mínimo de valoraciones temporalmente vinculadas, los autores lograron desanonimizar registros de usuarios concretos con una alta tasa de éxito. Este experimento reveló una verdad incómoda: la anonimización estadística que opera sobre datos reales mediante supresión o generalización es intrínsecamente frágil frente a la disponibilidad creciente de datos auxiliares externos. Cuanto mayor es la dimensionalidad del dataset original, mayor es el vector de ataque disponible para el adversario.

La confluencia de estas limitaciones fundamentales establece la motivación tecnológica central de este proyecto: la necesidad de un paradigma radicalmente diferente que no opere transformando datos reales existentes, sino generando datos sintéticos desde cero mediante modelos probabilísticos aprendidos. Esta distinción conceptual es la que otorga a la Inteligencia Artificial Generativa su superioridad estructural sobre los enfoques de anonimización clásica.

\subsection{Anonimización y el Dictamen del Grupo de Trabajo del Artículo 29}

El concepto de anonimización es frecuentemente malinterpretado en la industria tecnológica. Tradicionalmente, se ha creído que la supresión de identificadores directos (nombres, números de seguridad social) o la seudonimización mediante \textit{hashing} criptográfico bastaba para anonimizar una base de datos. La literatura de ciberseguridad ha refutado reiteradamente esta falacia, demostrando que atributos demográficos aparentemente inofensivos (como el código postal, el género y la fecha de nacimiento) son suficientes para reidentificar a un alto porcentaje de la población al cruzarlos con bases de datos externas.

Para arrojar luz legal sobre este conflicto tecnológico, el Grupo de Trabajo del Artículo 29 (el órgano consultivo independiente de la Unión Europea sobre protección de datos) emitió el Dictamen 05/2014 sobre técnicas de anonimización \parencite{article29wp2014}. En este documento fundacional, se dictamina que, para considerar un conjunto de datos como legítimamente anónimo, el proceso de transformación debe ser irreversible y garantizar la mitigación de tres riesgos fundamentales:
\begin{itemize}
    \item \textbf{Individualización (Singling out):} La imposibilidad de aislar a un individuo u obtener un registro único a partir de la masa de datos.
    \item \textbf{Vinculación (Linkability):} La incapacidad de relacionar dos o más registros relativos al mismo sujeto en bases de datos dispares.
    \item \textbf{Inferencia (Inference):} La inviabilidad de deducir, con una probabilidad significativa, el valor de un atributo sensible de un individuo a partir de la información proporcionada.
\end{itemize}
Cualquier metodología de síntesis de datos clínicos debe, por tanto, someterse a auditorías que demuestren el cumplimiento de estos tres criterios para obtener la inmunidad legal que otorga el GDPR.

\subsection{Privacidad Empírica vs. Privacidad Diferencial Formal}

Para garantizar la mitigación de los riesgos de inferencia expuestos, la comunidad científica ha desarrollado dos vertientes fundamentales de protección: la Privacidad Diferencial Formal y la Privacidad Empírica.

La Privacidad Diferencial (Differential Privacy o DP), propuesta por \textcite{dwork2006dp}, ofrece una garantía matemática inquebrantable. Un algoritmo se considera $\epsilon$-diferencialmente privado si la eliminación o adición de un único individuo en la base de datos de entrenamiento no altera la distribución de probabilidad de la salida del algoritmo más allá de un factor multiplicativo determinado por el presupuesto de privacidad $\epsilon$. Esta técnica se aplica inyectando ruido estadístico calibrado (mediante mecanismos de Laplace o Gaussianos) durante el entrenamiento de las redes generativas. Sin embargo, como apuntan \textcite{elemam2020synthetic}, forzar garantías teóricas rígidas en modelos de \textit{Deep Learning} genera una drástica degradación en la fidelidad y utilidad predictiva de los datos resultantes, arruinando su viabilidad corporativa.

Como alternativa, ha surgido el paradigma de la Privacidad Empírica \parencite{stadler2022synthetic}. Bajo este enfoque, no se inyecta ruido matemático a ciegas, sino que se entrena el modelo de IA a su máxima capacidad de fidelidad y, a posteriori, se ejecuta una batería de ciberataques simulados contra el modelo y sus datos resultantes. Si los atacantes algorítmicos fracasan en reidentificar a los pacientes, el riesgo se considera ``no razonablemente probable'', cumpliendo \textit{de facto} con el considerando 26 del GDPR sin sacrificar la calidad estadística.

\subsection{Modelos de Riesgo: Memorización e Inferencia de Membresía}

Dentro de la validación de la Privacidad Empírica, los modelos generativos son auditados frente a dos vulnerabilidades críticas. En primer lugar, la \textbf{Memorización}, cuantificada habitualmente mediante la métrica \textit{Distance to Closest Record} (DCR). Dado que las redes neuronales masivas tienen millones de parámetros, existe el riesgo inherente de que el modelo sobreajuste (overfitting) y literalmente ``memorice'' pacientes del set de entrenamiento, escupiéndolos de forma idéntica durante la fase de generación sintética. Un análisis DCR mide la distancia euclidiana entre el clon sintético más cercano y el paciente real original, debiendo garantizar que la distancia sea estrictamente superior a cero en todo el corpus.

En segundo lugar, se encuentra el \textbf{Ataque de Inferencia de Membresía (MIA)}. Formalizado matemáticamente por \textcite{shokri2017mia}, un MIA es un vector de ataque en el que un actor malicioso entrena un modelo clasificador binario en la sombra. El objetivo de este clasificador no es otro que predecir si un paciente específico en la vida real formó parte (membresía) del grupo de datos original con el que se entrenó la Inteligencia Artificial. Un modelo generativo robusto mitigará este ataque, forzando a la red del adversario a operar con una curva ROC-AUC cercana a 0.50, equivalente estadísticamente al lanzamiento de una moneda al azar.

\section{Evolución del Modelado Generativo Profundo}

La Inteligencia Artificial Generativa ha transitado por una evolución vertiginosa durante la última década. El objetivo fundamental de esta rama del Aprendizaje Automático no es discriminar o clasificar datos preexistentes, sino aprender la distribución de probabilidad subyacente $p_{data}(x)$ de un conjunto de datos reales, de modo que el modelo pueda generar nuevas muestras extraídas de una distribución aproximada $p_{model}(x)$.

\subsection{Redes Generativas Antagónicas (GANs)}

El punto de inflexión en la IA Generativa contemporánea aconteció en 2014 con la formulación de las Redes Generativas Antagónicas (GANs) por \textcite{goodfellow2014gan}. El marco de trabajo de las GANs se fundamenta en la Teoría de Juegos, proponiendo una arquitectura compuesta por dos redes neuronales que compiten en un juego minimax de suma cero: el Generador ($G$) y el Discriminador ($D$).

El Generador ingiere ruido aleatorio $z$ procedente de un espacio latente estructurado (usualmente una distribución gaussiana normalizada) y lo transforma en un vector $G(z)$ con el objetivo de imitar los datos reales. Por su parte, el Discriminador actúa como un clasificador binario que recibe una mezcla de datos reales $x$ y datos falsificados $G(z)$, con el mandato de predecir la autenticidad de cada muestra. La función de pérdida (loss function) que rige este enfrentamiento adversarial se define algebraicamente como:
$$ \min_G \max_D V(D, G) = \mathbf{E}_{x \sim p_{data}(x)}[\log D(x)] + \mathbf{E}_{z \sim p_z(z)}[\log (1 - D(G(z)))] $$

A medida que el entrenamiento avanza, el Generador se ve forzado a producir falsificaciones cada vez más indetectables, obligando al Discriminador a refinar sus heurísticas de detección. En el límite teórico, el Generador domina la distribución real y el Discriminador es incapaz de decidir ($D=0.5$).

\begin{figure}[htbp]
    \centering
    \includegraphics[width=0.85\textwidth]{GANs.png}
    \caption{Arquitectura general y flujo adversarial de una red GAN \parencite{gan_architecture_fig}.}
    \label{fig:gan_architecture}
\end{figure}

\subsection{Adaptaciones de las GANs al Dominio Tabular y sus Limitaciones}

A pesar de su impacto en la generación de imágenes y texto, la traslación de las GANs al dominio tabular médico presentó desafíos arquitectónicos significativos que motivaron una larga cadena de propuestas especializadas. A diferencia de los píxeles de una imagen, los datos tabulares combinan variables de naturaleza radicalmente heterogénea: valores continuos con distribuciones multimodales (como el número de procedimientos médicos), identificadores discretos ordinales (categorías de edad) y variables nominales de alta cardinalidad (como los 700 códigos diagnósticos ICD-9 originales).

Las primeras adaptaciones al dominio de la salud, como \textbf{MedGAN} \parencite{choi2017medgan}, emplearon autoencoders como preprocesadores para comprimir el espacio categórico antes de alimentar al Generador, mitigando parcialmente la incompatibilidad con el ruido gaussiano. Posteriormente, \textcite{xu2019ctgan} propusieron \textbf{CTGAN} (\textit{Conditional Tabular GAN}), que introdujo dos innovaciones decisivas: un preprocesamiento de Normalización de Modo-Específico mediante mezclas de Gaussianas (VGM) para modelar las distribuciones multimodales, y un mecanismo de muestreo condicional que obliga al Generador a producir ejemplos equilibrados de cada categoría, aliviando el problema del desbalanceo de clases. Propuestas más recientes como \textbf{CTAB-GAN} \parencite{zhao2021ctabgan} incorporaron capas convolucionales al Discriminador para capturar dependencias entre variables adyacentes al reorganizar el tensor tabular como una pseudo-imagen.

No obstante, todas las variantes GAN comparten dos limitaciones estructurales inherentes a la formulación minimax. La primera es la \textbf{inestabilidad del entrenamiento}: el juego adversarial puede colapsar si el Discriminador supera al Generador demasiado rápido, saturando el gradiente del Generador (el denominado \textit{vanishing gradient problem}). Para combatir esto, el presente trabajo empleó la variante WGAN-GP (\textit{Wasserstein GAN with Gradient Penalty}), que reemplaza la divergencia de Jensen-Shannon original por la Distancia de Wasserstein, proporcionando un gradiente de entrenamiento más suave y estable. La segunda limitación es el fenómeno del \textbf{colapso de modo} (\textit{mode collapse}): el Generador descubre un subespacio reducido de muestras que engañan consistentemente al Discriminador y colapsa en la repetición iterativa de esas mismas muestras, abandonando la representación de la diversidad estadística del dataset, incluyendo las clases minoritarias críticas.

\subsection{Autoencoders Variacionales (VAEs)}

Cofundados por \textcite{kingma2014vae}, los Autoencoders Variacionales (VAEs) ofrecen una perspectiva probabilística divergente, priorizando la estabilidad del entrenamiento mediante el modelado bayesiano. A diferencia del enfrentamiento minimax, un VAE consta de una red Codificadora (\textit{Encoder}) que comprime un registro multidimensional $x$ hacia una distribución estocástica en un espacio latente de baja dimensión $q_\phi(z|x)$. Posteriormente, una red Decodificadora (\textit{Decoder}) extrae un punto de esa distribución y procede a reconstruir el dato original $p_\theta(x|z)$.

El objetivo del entrenamiento en un VAE no es engañar a un discriminador, sino maximizar el Límite Inferior de Evidencia (\textit{Evidence Lower Bound} o ELBO). La función de pérdida se divide en dos componentes en perpetuo conflicto:
$$ \mathcal{L}_{VAE} = -\mathbf{E}_{q_\phi(z|x)}[\log p_\theta(x|z)] + D_{KL}(q_\phi(z|x) || p(z)) $$
El primer término actúa como una pérdida de reconstrucción clásica (fomentando copias precisas del paciente real), mientras que el segundo término es la Divergencia de Kullback-Leibler ($D_{KL}$), que actúa como un regularizador matemático forzando a las distribuciones latentes de los pacientes a asemejarse a una distribución normal estándar.

Esta tensión teórica (forzar a los datos reales a encajar en una única distribución continua regularizada) es, simultáneamente, la principal fortaleza y la mayor limitación de los VAEs en el dominio tabular. Mientras que en imágenes (rostros, objetos) la transición entre píxeles es continua, en los datos tabulares clínicos existe un desbalanceo agudo y una frontera discreta severa entre los pacientes sanos (clase mayoritaria) y los enfermos críticos (clase minoritaria). Al comprimir ambas realidades en un espacio latente unificado fuertemente regularizado, las redes VAE tienden a desdibujar o suavizar las fronteras de decisión de las minorías, degradando su rendimiento predictivo en métricas como el F1-Score.

\begin{figure}[htbp]
    \centering
    \includegraphics[width=0.85\textwidth]{VAEs.png}
    \caption{Estructura general y regularización del espacio latente en un Autoencoder Variacional (VAE) \parencite{vae_architecture_medium}.}
    \label{fig:vae_architecture}
\end{figure}

\section{Modelos Probabilísticos de Difusión (Denoising Diffusion)}

El panorama del aprendizaje profundo generativo experimentó un cambio de paradigma conceptual en el año 2020 con la publicación seminal de \textcite{ho2020ddpm}, la cual introdujo de manera práctica y escalable los Modelos Probabilísticos de Difusión de Eliminación de Ruido (Denoising Diffusion Probabilistic Models o DDPM). Estos modelos abandonan completamente la tensión adversarial de las GANs y la compresión restrictiva de los VAEs. Su formulación se inspira directamente en la termodinámica fuera del equilibrio de la física teórica.

\subsection{Fundamentos Matemáticos: Cadenas de Markov y Proceso Forward}

La fundamentación matemática de los modelos de difusión reside en dividir el complejo problema de generar datos a partir de ruido aleatorio en múltiples pasos discretos de magnitud reducida de corrección. El algoritmo define un Proceso Forward ($q$) o proceso de destrucción paramétrico de Markov. Dada una muestra real de la distribución original $x_0 \sim p_{data}(x)$, se le inyecta una varianza controlada y programada de ruido gaussiano en cada paso $t$, obteniendo secuencias latentes $x_1, x_2, ..., x_T$:
$$ q(x_t | x_{t-1}) = \mathcal{N}(x_t; \sqrt{1 - \beta_t} x_{t-1}, \beta_t \textbf{I}) $$
Donde $\beta_t$ es la varianza programada del ruido (el \textit{schedule}). Una propiedad fundamental de esta cadena de Markov es que permite computar $x_t$ directamente desde $x_0$ en un único paso analítico, sin necesidad de simular los $t$ pasos intermedios. Definiendo $\alpha_t = 1 - \beta_t$ y $\bar{\alpha}_t = \prod_{s=1}^{t} \alpha_s$, se obtiene:
$$ q(x_t | x_0) = \mathcal{N}(x_t; \sqrt{\bar{\alpha}_t} x_0, (1-\bar{\alpha}_t) \textbf{I}) $$
Esta simplificación, conocida como la \textit{reparametrización de Ho et al.}, permite muestrear $x_t$ en tiempo $O(1)$ como $x_t = \sqrt{\bar{\alpha}_t} x_0 + \sqrt{1-\bar{\alpha}_t} \varepsilon$, donde $\varepsilon \sim \mathcal{N}(0, \textbf{I})$. Si el número de pasos totales $T$ (típicamente 1000 pasos) es lo suficientemente grande, la distribución del tensor final $x_T$ se vuelve indistinguible de un ruido puramente aleatorio e isotrópico, condición matemáticamente necesaria para garantizar la generación desde ruido puro.

\subsection{El Proceso Inverso y el Aprendizaje de Denoising}

El corazón de la Inteligencia Artificial radica en aprender a ejecutar el Proceso Inverso ($p_\theta$). Si se es capaz de enseñar a una red neuronal profunda a descontaminar estadísticamente el tensor en el paso $x_t$ para predecir cómo era en el paso ligeramente anterior $x_{t-1}$, entonces la red puede tomar un ruido blanco completamente aleatorio ($x_T$) y reconstruir iterativamente una muestra realista ($x_0$):
$$ p_\theta(x_{t-1} | x_t) = \mathcal{N}(x_{t-1}; \mu_\theta(x_t, t), \Sigma_\theta(x_t, t)) $$
La red neuronal, referida como el \textit{Denoiser}, se entrena para estimar el ruido $\varepsilon_\theta(x_t, t)$ inyectado en el proceso forward, mediante el Error Cuadrático Medio:
$$ \mathcal{L}_{simple} = \mathbf{E}_{t, x_0, \varepsilon} \left[ ||\varepsilon - \varepsilon_\theta(x_t, t)||^2 \right] $$
Una vez entrenada la red, la media del paso inverso puede recuperarse analíticamente desde la predicción del ruido:
$$ \mu_\theta(x_t, t) = \frac{1}{\sqrt{\alpha_t}} \left( x_t - \frac{\beta_t}{\sqrt{1-\bar{\alpha}_t}} \varepsilon_\theta(x_t, t) \right) $$
Como demostró \textcite{nichol2021improved}, el rendimiento óptimo se alcanza al optimizar el planificador del ruido (\textit{schedule} coseno), y parametrizar la red de manera que reciba el paso temporal $t$ como una integración de codificación posicional (Timestep Embedding), permitiendo que la red sepa exactamente en qué estadio de degradación se encuentra la muestra.

\textbf{El Planificador de Ruido (\textit{Noise Schedule}).} La elección del planificador $\beta_t$ condiciona de forma decisiva la dinámica de aprendizaje del modelo. El planificador \textbf{lineal}, introducido por \textcite{ho2020ddpm}, incrementa $\beta_t$ de forma uniforme desde un valor mínimo $\beta_1$ hasta un máximo $\beta_T$, asegurando una destrucción progresiva y predecible de la información. El planificador de \textbf{coseno}, propuesto posteriormente por \textcite{nichol2021improved}, parametriza $\bar{\alpha}_t$ siguiendo una curva coseno que ralentiza la destrucción en los pasos iniciales, manteniendo más información durante el inicio del proceso y favoreciendo la reconstrucción de detalles finos en distribuciones de alta frecuencia, como las presentes en imágenes fotorrealistas. La elección entre ambas estrategias para el dominio tabular, y su interacción con la representación vectorial OHE de las variables categóricas, constituye uno de los ejes del estudio de ablación presentado en el Capítulo 5.

\begin{figure}[htbp]
    \centering
    \includegraphics[width=0.9\textwidth]{TABB.png}
    \caption{Procesos directo (Forward) de adición de ruido y reverso (Backward/Denoising) en modelos de difusión probabilísticos \parencite{diffusion_process_fig}.}
    \label{fig:diffusion_architecture}
\end{figure}

\subsection{Adaptación al Dominio Tabular: El Framework TabDDPM}

Históricamente, los DDPM se popularizaron de forma abrumadora en el procesamiento de imágenes, al ser respaldados internamente por arquitecturas U-Net. Sin embargo, en un desarrollo metodológico reciente, \textcite{kotelnikov2023tabddpm} publicaron TabDDPM, el primer acercamiento serio para adaptar el modelado probabilístico de la difusión a las complejas estructuras tabulares.

Dado que la difusión es un proceso estocástico que opera asumiendo espacios euclidianos y ruido gaussiano continuo, el desafío radica en el tratamiento de variables numéricas continuas/discretas y categóricas (códigos diagnósticos, razas, géneros). TabDDPM soluciona esto mediante un preprocesamiento continuo masivo: los valores numéricos son estandarizados mediante varianzas unitarias, y las variables categóricas son transmutadas a vectores difusos mediante el uso de la codificación densa iterativa (\textit{One-Hot Encoding}). Esta vasta matriz de variables mixtas es procesada por una red \textit{Multi-Layer Perceptron} (MLP) con alta capacidad con conexiones residuales similares a las arquitecturas \textit{Attention} de \textcite{vaswani2017attention}, logrando preservar las interdependencias evitando la degradación que experimentan las redes convolucionales al carecer el texto tabular de vecindad espacial.

Una ventaja arquitectónica fundamental de TabDDPM frente a las GANs es que su proceso de entrenamiento es intrínsecamente estable: al no existir competición adversarial, el gradiente de error es un simple MSE sobre el ruido, eliminando los problemas de colapso de modo y desvanecimiento del gradiente que afectan sistemáticamente a las variantes GAN en datasets con distribuciones multimodales y clases desbalanceadas.

\section{El Ecosistema de Evaluación de Datos Tabulares Sintéticos}

Al carecer del componente perceptual del cual gozan las imágenes sintéticas (donde un humano puede evaluar si un rostro o paisaje es ``realista'' mediante inspección visual), la auditoría algorítmica de los gemelos tabulares requiere una metodología de validación estrictamente cuantitativa. Esta metodología se articula en torno a la plataforma \textbf{SDV} (\textit{Synthetic Data Vault}) \parencite{sdv2023}, un ecosistema de código abierto que proporciona tanto los motores generativos (CTGAN, TVAE) como los módulos de evaluación de calidad (\textit{sdmetrics}) utilizados en el presente proyecto para computar las métricas de correlación y diversidad. El ecosistema de evaluación se divide en tres ejes cardinales.

\subsection{Fidelidad Estadística: La Distancia de Wasserstein}

Para evaluar si el modelo ha aprendido con precisión las distribuciones del mundo real, se requiere medir la divergencia entre la distribución de densidad de los datos reales $p(x)$ y la distribución de los datos sintéticos $q(x)$. La métrica estándar de facto para variables continuas en el Deep Learning moderno es la Distancia de Wasserstein de Primer Orden ($W_1$), originada en los problemas matemáticos del transporte óptimo (Earth Mover's Distance) \parencite{villani2008wasserstein}.

Conceptualmente, la distancia de Wasserstein calcula la distancia de transporte óptimo mínima indispensable para reformar una distribución de probabilidad y transmutarla para que iguale a la otra. Cuanto menor sea el coste (acercándose asintóticamente a 0.0), más indistinguibles estadísticamente son la muestra generada y la realidad subyacente. Adicionalmente, el análisis de fidelidad debe ser robustecido a nivel multivariante, calculando la Diferencia de la Matriz de Correlación MAE (\textit{Mean Absolute Error}). Esta métrica resta iterativamente la matriz de correlación de Pearson o Cramér del dataset original a la matriz generada; una pérdida del orden del 1\% al 3\% es indicativo de un modelo eficaz capaz de preservar las reglas de negocio lógicas inherentes.

\subsection{Utilidad Predictiva: El Paradigma TSTR frente a TRTR}

El grado de madurez absoluto de un framework generativo no reside en su capacidad estadística descriptiva, sino en su utilidad algorítmica operativa. Para este fin, la validación se somete al riguroso protocolo formal denominado \textbf{TSTR} (\textit{Train on Synthetic, Test on Real}). En este paradigma, es fundamental distinguirlo de sus variantes complementarias dentro de los protocolos de evaluación estándar de la literatura de síntesis de datos:

\begin{itemize}
    \item \textbf{TRTR} (\textit{Train on Real, Test on Real}): El clasificador externo se entrena y evalúa sobre datos reales. Constituye el \textit{baseline} de referencia o ``límite superior de rendimiento'' alcanzable, ya que maximiza la disponibilidad de señal de entrenamiento.
    \item \textbf{TSTR} (\textit{Train on Synthetic, Test on Real}): El clasificador se entrena exclusivamente con datos sintéticos y se evalúa sobre datos reales no observados. Es el protocolo estándar para certificar que el gemelo sintético preserva suficiente señal predictiva como para reemplazar al dataset real en un pipeline de ML.
    \item \textbf{TSTS} (\textit{Train on Synthetic, Test on Synthetic}): Ambas particiones son sintéticas. Esta variante no mide la utilidad en el mundo real; solo mide la consistencia interna del modelo generativo.
\end{itemize}

En el protocolo TSTR, la IA generativa despliega el volumen total de gemelos virtuales sin incluir un solo paciente biológico genuino. Este cuerpo sintético se emplea de forma hermética para entrenar un clasificador predictivo externo. Dada su contrastada resistencia al sobreajuste, la herramienta predictiva predilecta suele ser \textit{XGBoost} \parencite{chen2016xgboost}. Posteriormente, este modelo se despliega para ejecutar su función predictiva sobre el grupo de datos reales de prueba (\textit{Holdout Test Set}) extraído al inicio del experimento y que jamás fue observado. El desempeño del clasificador entrenado por la inteligencia sintética se compara directamente contra el \textit{baseline} (un modelo idéntico instruido con pacientes biológicos genuinos). Si el margen de pérdida predictiva, tasado en métricas de clasificación pesada como la curva ROC-AUC o el F1-Score armónico, es inferior al 5.0\%, el ecosistema generativo se certifica como un repuesto integralmente compatible (\textit{Drop-In Replacement}) con garantías operativas óptimas en ambientes clínicos industriales.

\subsection{Privacidad Empírica: DCR y MIA como Vectores de Auditoría}

El tercer eje de la evaluación, y el más crítico desde la perspectiva regulatoria, es la auditoría de la privacidad empírica. Como se detalló en la sección 1.4, los dos vectores de ataque fundamentales que deben neutralizarse son la Memorización (DCR) y la Inferencia de Membresía (MIA). Un modelo generativo que supere favorablemente ambas auditorías —mostrando un DCR mínimo estrictamente mayor que cero y un AUC del atacante MIA estadísticamente indistinguible de 0.50— habrá demostrado que sus datos sintéticos cumplen los tres criterios del Dictamen del Artículo 29 (individualización, vinculación e inferencia), obteniendo el aval legal para su clasificación como datos anónimos fuera del alcance del GDPR.
